\section{Introduction}
Learning robust models that generalize well is critical for many real-world applications \cite{Zech2018,degrave2021ai}.
Yet, the classical Empirical Risk Minimization (ERM) lacks robustness to distribution shifts \cite{hendrycks2018benchmarking,NEURIPS2020_6cfe0e61,d2020underspecification}.
To improve out-of-distribution (\ood) generalization in classification, several recent works proposed to train models simultaneously on multiple related but different domains \cite{muandet2013domain}.
Though theoretically appealing, domain-invariant approaches \cite{peters2016causal} either underperform \cite{arjovsky2019invariant,krueger2020utofdistribution} or only slightly improve \cite{coral216aaai,rame_fishr_2021} ERM on the reference DomainBed benchmark \cite{gulrajani2021in}.
The state-of-the-art strategy on DomainBed is currently to average the weights obtained along a training trajectory \cite{izmailov2018}.
\cite{cha2021wad} argues that this weight averaging (WA) succeeds in \ood because it finds solutions with flatter loss landscapes.

In this paper, we show the limitations of this flatness-based analysis and provide a new explanation for the success of WA in \ood.
It is based on WA's similarity with ensembling \cite{Lakshminarayanan2017}, a well-known strategy to improve robustness \cite{Ovadia2019,ashukha2020pitfalls}, that averages the predictions from various models.
Based on \cite{ueda1996generalization}, we present %in \Cref{subsec:biasvariance}
a bias-variance-covariance-locality decomposition of WA's expected error.
It contains four terms:
\textit{first} the bias that we show increases under shift in label posterior distributions (\ie correlation shift \cite{ye2021odbench});
\textit{second}, the variance that we show increases under shift in input marginal distributions (\ie diversity shift \cite{ye2021odbench});
\textit{third}, the covariance that decreases when models are diverse;
\textit{finally}, a locality condition on the weights of averaged models.
% to enforce the validity of our approximation.

Based on this analysis, we aim at obtaining diverse models whose weights are averageable with our Diverse Weight Averaging (DiWA) approach.
In practice, DiWA averages in weights the models obtained from independent training runs that share the same initialization. The motivation is that those models are more diverse than those obtained along a single run \cite{Fort2019DeepEA,gontijolopes2022no}.
Yet, averaging the weights of independently trained networks with batch normalization \cite{batchnorm2015} and ReLU layers \cite{agarap2018deep} may be counter-intuitive.
Such averaging is efficient especially when models can be connected linearly in the weight space via a low loss path. Interestingly, this linear mode connectivity property \cite{Frankle2020} was empirically validated when the runs start from a shared pretrained initialization \cite{Neyshabur2020}.
This insight is at the heart of DiWA but also of other recent works \cite{Wortsman2022robust,mergefisher21,Wortsman2022ModelSA}, as discussed in \Cref{sec:related_work}.

In summary, our main contributions are the following:
\begin{itemize}
    \item We propose a new theoretical analysis of WA for \ood based on a bias-variance-covariance-locality decomposition of its expected error (\Cref{sec:theory}). By relating correlation shift to its bias and diversity shift to its variance, we show that WA succeeds under diversity shift.
    \item We empirically tackle the covariance term by increasing the diversity across models averaged in weights. In our DiWA approach, we decorrelate their training procedures: in practice, these models are obtained from independent runs (\Cref{sect:dwa}).
    We then empirically validate that diversity improves \ood performance (\Cref{sect:analysis}) and show that DiWA is state of the art on all real-world datasets from the DomainBed benchmark \cite{gulrajani2021in} (\Cref{sec:domainbed}).%
\end{itemize}%